\PassOptionsToPackage{table}{xcolor}
\documentclass{article}
\usepackage{iclr2027_conference,times}
\usepackage{hyperref}
\usepackage{url}
\usepackage{graphicx}
\usepackage{booktabs}
\usepackage{amsmath}
\usepackage{xcolor}
\definecolor{linkblue}{HTML}{1F5A94}
\hypersetup{
  colorlinks=true,
  linkcolor=linkblue,
  citecolor=linkblue,
  urlcolor=linkblue
}
\usepackage{amsfonts}
\usepackage{amssymb}
\usepackage{amsthm}
\usepackage{enumitem}
\usepackage{pifont}
\DeclareMathSizes{9}{10}{7}{5}

% Theorem environments and shared math notation.
\newtheorem{theorem}{Theorem}
\newtheorem{assumption}{Assumption}
\newtheorem{proposition}{Proposition}
\theoremstyle{definition}
\newtheorem{definition}{Definition}
\theoremstyle{plain}
\newcommand{\E}{\mathbb{E}}
\newcommand{\Var}{\mathrm{Var}}
\newcommand{\norm}[1]{\lVert #1 \rVert}
\newcommand{\sigci}[1]{$\mathbf{#1}$}
\newcommand{\sigrow}[5]{\textbf{#1} & \textbf{#2} & \sigci{#3} & \textbf{#4} & \textbf{#5} \\}

% Tighten float/caption spacing to fit the 9-page limit (formatting only).
\setlength{\textfloatsep}{8pt plus 2pt minus 2pt}
\setlength{\floatsep}{8pt plus 2pt minus 2pt}
\setlength{\intextsep}{8pt plus 2pt minus 2pt}
\setlength{\abovecaptionskip}{4pt}
\setlength{\belowcaptionskip}{2pt}

% Compact captions: keep plots readable while avoiding body-sized caption blocks.
\makeatletter
\long\def\@makecaption#1#2{%
  \vskip\abovecaptionskip
  {\footnotesize
    \sbox\@tempboxa{\textbf{#1:} #2}%
    \ifdim \wd\@tempboxa >\hsize
      \textbf{#1:} #2\par
    \else
      \global\@minipagefalse
      \hb@xt@\hsize{\hfil\box\@tempboxa\hfil}%
    \fi}
  \vskip\belowcaptionskip}
\makeatother

% -----------------------------------------------------------------------
% Working title — change before submission
% -----------------------------------------------------------------------
\title{Weights Have Traits: Estimating Fine-Tuning Lineage in LLMs}

% Authors hidden for double-blind review.
% \iclrfinalcopy  % Uncomment for camera-ready only.

\newcommand{\fix}{\marginpar{FIX}}
\newcommand{\new}{\marginpar{NEW}}

% Placeholder macro for unwritten sections — shows in red, visible in draft
\newcommand{\TODO}[1]{\textcolor{red}{\textbf{[TODO: #1]}}}

% Shorthand for figure path prefix
\newcommand{\figpath}{figures}

\begin{document}
\maketitle

% -----------------------------------------------------------------------
\begin{abstract}
While over a million deep learning models are hosted online, most are fine-tuned descendants from a small set of ancestors.  The lineage of these models is only sometimes documented fully, even though provenance, auditing, and safe deployment all depend on proper documentation. In response to this issue, we hypothesize that fine-tuning acts as an evolutionary process that leaves signal in weight space, and show that a model \textit{phylogeny} (evolutionary tree) can be estimated from model weights alone and scored against a known ground-truth phylogeny, recovering all true clades (branches) in 85--100\% of runs.  We treat individual weight layers as gene-like estimators of lineage, allowing each layer to produce its own phylogeny.  We show that the phylogeny can be estimated accurately from a tiny fraction of the weights (0.1--0.2\% of the weights in our experiments) can recover all true clades, suggesting that lineage can be estimated from a small, architecture-specific weight subset. Moreover, we show that this phylogeny estimation is associated with known mathematical concepts from phylogenetics: the Atteson margin and four-point additivity. Beyond phylogeny estimation, we report two additional empirical findings. First, we show weight-space proximity predicts behavioral similarity on several held-out tasks involving truthfulness, knowledge, reasoning, and common sense, though the effect size depends on the task. Second, we show weight-based phylogeny estimation can match or outperform behavior-based estimation in both experimental and real scenarios. On the matched controlled trees, weights recover 100\% of clades versus PhyloLM's 35\% under rank-8 LoRA and 100\% versus 52\% under full fine-tuning, while both methods recover all three documented splits in the real Mistral family.

\end{abstract}

% -----------------------------------------------------------------------
\section{Introduction}
\label{sec:intro}




While there are over one million models available on HuggingFace for deep learning and a handful of frontier
labs constantly releasing updated versions, most models are
fine-tuned descendants of a few ancestors.  Very few of those models transparently document their full model lineage \citep{horwitz2025chartatlasworldsmodels}. Benchmarks can reveal how a model performs on a task, but not how a model originated, how it differs from its sibling models, or how it may
fail. In contrast, training lineage can reveal much about model provenance, auditing, and safety~\citep{behavioral2026provenance, nikolic2026model, tensorguard2025}.

We, like other prior works~\citep{horwitz2025mother, yu2025neural, wu2026llmdna,gallagher2026}, observe that fine-tuning is a path-dependent, branching process. Each
child model inherits its parent's  updated weights, nudged towards a certain loss and/or data set. This is
precisely an evolutionary process, as described in \cite{yax2025phylolm}, and moreover, the evolutionary path leaves signal in weight space. We treat a line or genealogy of fine-tuned models
as a phylogeny (i.e. an evolutionary tree~\cite{saitou1987neighbor}) and ask whether this phylogeny can be \emph{estimated} from the model weights alone.

As motivation, consider a firm that forks a base model into finance, IT, and fraud variants, each fine-tuned until
dozens of versions circulate with no recorded lineage (Fig.~\ref{fig:overview}A). A practitioner needs
to know which to build on, where two versions diverged, and where each fails. Phylogenetic algorithms
were built for exactly this inference in biology~\citep{saitou1987neighbor, robinson1981comparison},
and we adapt them to neural net model weights. We estimate a phylogeny from model weights and score it against a \emph{known}
ground-truth phylogeny (Fig.~\ref{fig:overview}A), link the changes in weights to changes in behavior (Fig.~\ref{fig:overview}B), and ask whether diagnostic statistics indicate when phylogeny estimation will succeed even without access to the true tree (Fig.~\ref{fig:overview}C).

\begin{figure}[t]
  \begin{center}
    \includegraphics[width=\linewidth]{\figpath/fig1_overview_v6.pdf}
  \end{center}
  \caption{Overview. \textbf{(A)} Can the training tree be recovered from the leaf models' weights alone, with ancestors latent? \textbf{(B)} Do models closer in weight space behave alike on held-out tasks? \textbf{(C)} Does label-free tree-likeness signal when recovery works?}
  \label{fig:overview}
\end{figure}



Our primary contribution is that analyzing weight changes between fine-tunings allows accurate phylogeny estimation with latent ancestors. We use a fitted Atteson diagnostic and label-free four-point additivity as empirical diagnostics to characterize when phylogeny estimation succeeds. As a result of these weight-based estimates, we can further show two behavior-related consequences. In summary we show:
\begin{itemize}[leftmargin=1.3em,itemsep=1pt,topsep=2pt,parsep=0pt]
  \item \textbf{Phylogeny estimation from leaf model weights against ground truth.} Guided by the geometry of
    fine-tuning updates, our method of phylogeny estimation utilizes the  Fr\'echet mean of trees to estimate the model phylogeny from weights alone. Our controlled experiments recover \textit{all} true clades (branches) in $85$--$100\%$ of runs
    ($92$--$100\%$ clade recovery). We evaluate against known ground-truth phylogenies and show that a single attention key matrix, about 0.1--0.2\% of the scalar weights, can recover the known phylogenies. We also introduce the concept of using weight layers as gene-like estimators which yields a separate phylogeny from each layer and  highlights where lineage signal is strongest. We show that our weight-based method accurately recovers all the splits in a small but real Mistral phylogeny.
  \item \textbf{Scenarios where weight-based estimation outperforms behavior-based estimation.}
    On the same 46 controlled trees, weights recover 100\% of clades versus 35\% for
PhyloLM~\citep{yax2025phylolm} under rank-8 LoRA, and 100\% versus 52\% under full fine-tuning.  Moreover, we show that the Atteson diagnostic and label-free diagnostic also signal this association.
    On a documented Mistral family with large RLHF/DPO forks, both methods
    recover all three true splits. Thus behavior-only estimation can succeed after behaviorally large
    updates, while weight access remains substantially more informative for smaller updates and
    requires no output sampling.
  \item \textbf{Evidence of similar weights signaling similar behavior.} Changes in weights can be used to predict similarity in model behavior. In particular, we show that models that are closer in weight-space behave more similarly in behavior-space on \textit{held-out} tasks such as truthfulness, reasoning, knowledge, and common sense.
\end{itemize}
% Submission TODO: add the anonymous code-and-artifact URL here once the blind snapshot exists.
Our main controlled topology results cover Flan-T5-base (250M) and Llama-3.2-1B, with both full
fine-tuning and variants of LoRA. 






% -----------------------------------------------------------------------
% Swati Rallapalli's related-work draft; keep this as a standalone source.
\section{Related Work}
\label{sec:related}

Deep learning lineage studies broadly infer model relationships from either \emph{model weights} or \emph{behavior}. These methods have been driven by different applications including provenance and IP \citep{horwitz2025mother, zeng2024huref, nikolic2026model}, model poisoning and misalignment \citep{cloud2025subliminal}, trait inheritance \citep{cloud2025subliminal, agents2026same}, and data auditing \citep{zhu2025independence}. 

\textbf{Model weights:} Several works trace weight changes from fine-tuning or reinforcement learning from human feedback (RLHF) back to parent models. MoTHer \citep{horwitz2025mother} reconstructs directed model trees from weight similarity via a directed minimum spanning tree and evaluates against controlled, known model hierarchies; \citet{yu2025neural} identify parent--child and grandparent--child lineage with both a learning-free fine-tuning approximation and a learned detector ($>99\%$ accuracy). Beyond direct lineage, spectral weight ``fingerprints'' identify common ancestry \citep{ghostspec2025}, and Spectral DeTuning \citep{horwitz2024spectral} reconstructs base weights from shared LoRA variants. Unlike MoTHer, which reconstructs relationships among observed parent and descendant models, we observe only leaf models and infer a phylogeny with latent ancestral branching points. This setting is relevant to LLM family reconstruction, where intermediate or parent models may not be available.

REEF \citep{zhang2025reef} compares internal representations using centered kernel alignment (CKA; \citealp{kornblith2019similarity}) to identify model derivation, while TensorGuard \citep{tensorguard2025} uses gradient-response fingerprints for pairwise model similarity and family classification.

Other works examine \emph{why} lineage is detectable in models. Models fine-tuned from a shared initialization stay within a single low-error basin and are robust to SGD noise \citep{neyshabur2020,wortsman2022model,frankle2020linear}, so pairwise weight distance can approximate genealogical distance. Additionally, task vectors \citep{ilharco2023editing} are near-orthogonal across tasks, so distinct lineages contribute independent directions. LoRA \citep{hu2022lora} and intrinsic-dimensionality results \citep{aghajanyan2021intrinsic} demonstrate that meaningful adaptation can occur in low-dimensional parameter subspaces. Together, these results motivate the possibility that lineage information remains detectable after fine-tuning.

\textbf{Behavior:} Behavioral approaches infer model relationships using only API access or model outputs, without requiring access to model weights \citep{yax2025phylolm,wu2026llmdna}. PhyloLM \citep{yax2025phylolm} builds LLM phylogenies from output similarity and predicts benchmark performance, while LLM DNA \citep{wu2026llmdna} defines a functional ``DNA'' representation carrying lineage information. Both methods estimate phylogenies over hundreds of models. Other output-only methods test pairwise provenance from next-token distributions \citep{nikolic2026model} or output independence \citep{zhu2025independence}. A key advantage of behavioral methods is that they can compare models across architectures and sizes, whereas weight-based comparisons generally require compatible parameterizations. Conversely, weight-based lineage estimates do not depend on the choice of prompts or test tasks.

\textbf{Our Approach:} We compare weight- and behavior-based approaches, determine where phylogenetic signal can be detected, and evaluate the estimated phylogenies against known true trees. Our findings are consistent with prior work showing that inherited behavioral signals can be preserved under sufficient training signal \citep{cloud2025subliminal}, as well as phylogenetic results showing that the underlying tree can be correctly estimated when certain conditions on the estimated distances are satisfied \citep{atteson1999performance}.  As Table~\ref{tab:related} shows, ours is the only work that combines weight-based estimation, latent ancestors (which allow for unobserved intermediate models), controlled LLM fine-tuning trees, a link between weights and behavior, and a label-free tree-likeness score that serves as an empirical measure of confidence in the estimated phylogeny.

\begin{table}[t]
\caption{We compare with the closest prior methods. Controlled trees are families the authors fine-tuned. Empirical confidence is a label-free score that is tested against estimation accuracy.}
\label{tab:related}
\begin{center}
\small
\setlength{\tabcolsep}{4pt}
\begin{tabular}{lcccccc}
\toprule
Method & Access & \shortstack{Latent\\ancestors} & \shortstack{Controlled\\trees} & \shortstack{Weight--\\behavior} & \shortstack{Empirical\\confidence} & \shortstack{Cross-\\arch.} \\
\midrule
MoTHer \citep{horwitz2025mother}
  & weights  & \ding{55}  & vision     & \ding{55}  & \ding{55}  & \ding{55} \\
Neural Phylogeny \citep{yu2025neural}
  & weights  & \ding{55}  & vision     & \ding{55}  & \ding{55}  & \ding{55} \\
PhyloLM \citep{yax2025phylolm}
  & behavior & \checkmark & \ding{55}  & \ding{55}  & \ding{55}  & \checkmark \\
LLM DNA \citep{wu2026llmdna}
  & behavior & \checkmark & \ding{55}  & \ding{55}  & \ding{55}  & \checkmark \\
\textbf{Ours}
  & \textbf{weights} & \checkmark & \textbf{LLM} & \checkmark & \checkmark & \ding{55} \\
\bottomrule
\end{tabular}
\end{center}
\end{table}


% -----------------------------------------------------------------------x
\section{Framework}
\label{sec:framework}
% -----------------------------------------------------------------------

We first define our terms, notation, and statistical framework. A \emph{phylogeny} $\mathcal{P}$ is a tree
whose leaves are the observed models and whose internal nodes are latent common ancestors; a
\emph{clade} is the set of leaves descending from one such ancestor (one branch of the tree). We analyze
single-root model families but observe only their terminal models. Thus,  we estimate the unrooted leaf topology with internal ancestors latent. We form a pairwise distance matrix and apply neighbor-joining (NJ)
\citep{saitou1987neighbor} to estimate the phylogeny.

\textbf{Phylogeny estimation.}\label{sec:framework:phylo}
Assume we have a model $g(\theta)$ where $\theta$ is composed of multiple layers $\theta_\ell$. To estimate $\widehat{\mathcal{P}}$: (1) choose a set of layers $\theta_\mathcal{L}$ and a distance metric $d$;
(2) for each layer $\ell$ compute $D_{i,j}^{(\ell)} = d(\theta^{(\ell)}_i, \theta^{(\ell)}_j)$ between models
$i,j$; and (3) apply NJ to each selected layer's matrix to obtain a layer-level tree. We either
analyze these layer-level trees separately or combine them in tree space, as described below. Analyses
that require one pairwise weight distance instead use the arithmetic mean of the active layer-level
distance matrices.

Step (3) lets different layers yield different estimates of $\widehat{\mathcal{P}}$, which we use to
assess each layer's lineage signal. We compare estimated and true phylogenies using Robinson--Foulds (RF) distance \citep{robinson1981comparison}, which counts disagreeing bipartitions as false positives (FP) and false negatives (FN). We restrict our analysis to unrooted phylogenies with at least four leaves, since three leaves have only one possible unrooted topology and therefore no nontrivial internal split to recover. Because NJ returns binary trees while our ground-truth phylogenies may contain \textit{polytomies} (nodes with more than two children), NJ must introduce additional splits to resolve them. RF counts a split present in the estimated tree but absent from the ground-truth tree as a false positive (FP); thus, splits introduced solely to resolve a polytomy may be counted as FPs even though they do not contradict the known lineage. We therefore focus on false negatives (FN), which count ground-truth clades missed by the reconstruction. We also report clade recovery, $1-\mathrm{FN}/|\text{true splits}|$, the fraction of ground-truth clades recovered, and polytomy-aware exact recovery (PAER), the fraction of runs in which all ground-truth clades are recovered ($\mathrm{FN}=0$). Separately, we report the standardized within-run branch correlation $r_{\mathrm{branch}}$ between weight distance and a same-branch indicator to assess whether weight distance reflects lineage relationships.



\textbf{Distance metric.}\label{sec:framework:metric}
We flatten each weight layer to a vector and compare models with vector distance and adopt cosine distance ($1$ minus cosine
similarity) throughout, for three reasons. It is bounded and scale-invariant across layers, so no single
high-magnitude layer dominates the aggregate, unlike $\ell_1$ or $\ell_2$, the latter which is used in \cite{horwitz2025mother}. It is the prevailing choice in
recent provenance and fingerprinting work \citep{zeng2024huref, behavioral2026provenance,
zhu2025independence}, keeping our distances comparable to that literature. Finally, cosine distance has a simple geometric interpretation: descendant models whose weight vectors remain similarly oriented in parameter space are considered closer. For LoRA-trained trees, we compare adapter-merged weights as opposed to the adapters to one another.



\textbf{Per-layer estimated phylogenies.}
We treat each weight layer as an independent estimator of lineage, analogous to using individual
genes to estimate biological phylogenies~\citep{maddison1997gene}. For a given training tree and
layer, we compute the cosine distance matrix across its observed models and apply NJ, producing one
estimated phylogeny for that layer.  The per-layer profile (see App.~\ref{app:layertrace}) computes the scores for each tree and averages them across simulations. It never averages distance
matrices across layers or combines different layers' trees. We exclude layers for which all pairwise
cosine distances are zero, since they provide no lineage information under this metric.

When one aggregate topology is required, we instead compute the Billera--Holmes--Vogtmann (BHV)
Fr\'echet mean of the layer-level trees~\citep{billera2001geometry, willis2019confidence}, a tree specific metric, considering only non-frozen layers in the case of LoRA.   For analyses that require a single pairwise
weight-distance matrix, including shared lineage, four-point additivity, the Atteson diagnostic,
and weight--behavior associations, we use the mean of the nonzero active-layer distance
matrices. Behavioral similarity is likewise the mean of prompt-paired output-embedding
cosine similarities. Thus, the per-layer profile, aggregate topology, and pairwise behavioral
analysis use separate, explicit aggregation rules.



\textbf{When estimation succeeds.}\label{sec:framework:theory}
We assess lineage signal at two levels. First, we test whether models on the same branch are closer in weight space than models on different branches. Within each valid run, we compare these pairwise distances using the Mann--Whitney $U$ statistic. This provides a simple measure of pairwise lineage signal without requiring estimation of the full phylogeny.

Second, we ask when the full phylogeny can be estimated from the weight distances. The NJ method is exact when the distance matrix is
additive, a condition model-weight distances satisfy only approximately. Atteson's theorem makes the
condition precise~\citep{atteson1999performance}. Let $D_T$ denote the additive leaf-distance matrix of a binary tree $T$ with positive edge lengths ${L_s}$, and let $D$ be an observed distance matrix. If
$$
\lVert D-D_T\rVert_\infty < \frac{1}{2}\min_s L_s,
$$
then NJ applied to $D$ estimates the topology of $T$.

Equivalently, define
$$
m:=\frac{\min_s L_s}{2\lVert D-D_T\rVert_\infty}.
$$
Then $m>1$ is sufficient for exact recovery. This condition is conservative as NJ may still recover the correct topology when $m\leq1$.

In our controlled experiments the true topology is
known, but its edge lengths are not. We fit nonnegative edge lengths on that topology to obtain
$D_{\mathrm{fit}}$ and report the analogous diagnostic
$\widetilde m=\min_s\widehat L_s/(2\lVert D-D_{\mathrm{fit}}\rVert_\infty)$. Because
$D_{\mathrm{fit}}$ is estimated from $D$ and some truth trees contain polytomies, $\widetilde m$ is an approximation rather than the true Atteson condition.  Because of this fitting approximation, we also
adopt a label-free proxy computed from $D$ alone, based on the classical four-point condition for tree
metrics~\citep{sempleSteel2003phylogenetics}.

\textit{Four-point additivity} ($A$). Non-additivity of a distance matrix is measured by the quartet
$\delta$-score of \citet{holland2002delta}. For each quartet of leaves the three pairwise-distance sums,
sorted $s_1\ge s_2\ge s_3$, give
\[
  \delta_q=\frac{s_1-s_2}{s_1-s_3}\in[0,1] ,
\]
which is $0$ exactly when the four-point condition $s_1=s_2$ holds \citep{sempleSteel2003phylogenetics}
(the additive, tree-metric case) and $1$ at the maximally non-treelike quartet (the two \emph{smallest}
sums equal, the largest strictly greater). We report its complement, averaged over the $Q$ quartets,
\[
  A=\frac{1}{Q}\sum_{q=1}^{Q}\bigl(1-\delta_q\bigr)
   =\frac{1}{Q}\sum_{q=1}^{Q}\frac{s_2^{(q)}-s_3^{(q)}}{s_1^{(q)}-s_3^{(q)}}\in[0,1] ,
\]
so that $A=1$ is perfectly additive and larger $A$ is associated with more accurate phylogeny estimation. Because $A$ requires no labels,
it serves as our proxy when the true tree is unknown. For evaluation against a known tree we
additionally report $A$ over only the \emph{informative} quartets (those the truth resolves into two
cherries), since quartets that fall inside a polytomy carry no split to recover and only dilute the
all-quartet mean. We relate $A$ and the fitted Atteson diagnostic to phylogeny estimation accuracy in \S\ref{sec:results}.


% -----------------------------------------------------------------------
\section{Experiments}
\label{sec:experiments}
% -----------------------------------------------------------------------
We structure our experiments around three questions: whether the model lineage or phylogeny can be
estimated from weight changes; whether those changes predict downstream held-out behavior; and
how behavior-based phylogeny estimation differs from weight-based.

\paragraph{Controlled and real training trees.}
\label{sec:exp:trees}
% Base: Flan-T5-base (250M). Bernoulli branching, 4 tasks, 50 replications,
% 5--8 leaves, known topology by construction.
% 2000 steps, lr=3e-4, batch=8, grad_accum=4, bf16.
% Controlled design enables exact evaluation against the latent-ancestor leaf topology.

To create a controlled experiment, we generate stochastic model training trees beginning from a root
model (process detailed in App.~\ref{app:details}). The topology-recovery cohort contains the 46 trees
with a nontrivial unrooted truth split, each with 4--10 observed leaves; the behavioral analysis uses
all 50 generated trees. We randomly assign a task/dataset pair to each edge, and at each split
fine-tune a child model from its parent on that pair. LoRA chains are sequentially merged so the
primary LoRA leaves are standalone full-weight checkpoints; the separately marked all-projections
analysis is analyzed in cumulative update space. We then estimate a phylogeny from the \textit{model
leaves} and compare it to the known training tree, testing both shared lineage and phylogeny
estimation.

We vary the root model, task types, and trained layers. Our primary evaluation uses an
encoder--decoder model (Flan-T5-base, 250M) and a decoder-only model (Llama-3.2-1B), fine-tuning on
Q\&A, classification, translation, and summarization. For Flan-T5, we run full fine-tuning and four LoRA \citep{hu2022lora}
scopes (key only; query/key/value; query/key/value/output; all attention and feed-forward projections).
For Llama-3.2-1B, we run full fine-tuning and query/key/value LoRA at ranks 8 and 64. For each base
model, we use the same task instructions. We also
validate our method on three documented model lineages from HuggingFace. We summarize them in the
main text and show all three in App.~\ref{app:organic-additional}.

\paragraph{Shared lineage.}
To determine shared lineage, MoTHer~\citep{horwitz2025mother} assembles pairwise relatedness and direction scores into a directed
model tree, while Neural Phylogeny~\citep{yu2025neural} identifies parent--child relationships. Our
auxiliary ordering analysis instead works at the branch level among observed leaves. Within each
training tree, two leaves are labeled ``same branch'' when their first ancestor below the root is
the same; all other leaf pairs are labeled ``cross branch.'' We test whether cross-branch distances
are larger with a one-sided Mann--Whitney $U$ test and report the corresponding rank-biserial effect
$r_{\mathrm{rb}}=2U/(n_{\mathrm{cross}}n_{\mathrm{same}})-1$. We also report the point-biserial correlation. Across trees, we report the mean and
sample SE of $r_{\mathrm{rb}}$, and Fisher $z$-average $r_{\mathrm{branch}}$\footnote{For each run we
compute $z=\operatorname{arctanh}(r)$, average on the $z$ scale, and back-transform with
$r=\tanh(\bar z)$.}.

\paragraph{Phylogeny estimation.}
\label{sec:exp:weights}
Besides shared lineage, we test whether our method estimates the phylogeny. From \S\ref{sec:framework:theory}, accurate estimation should be possible when the distance matrix is sufficiently additive, for which the label-free four-point additivity $A$ is a practical proxy. For each tree we report RF, FN, clade recovery, and polytomy-aware exact recovery (PAER; $\mathrm{FN}=0$), relating them to the fitted Atteson diagnostic $\widetilde m$ and $A$ in \S\ref{sec:results}. The per-layer trees also let us identify which layers are strong or weak carriers of lineage signal, via the highest- and lowest-FN layers per architecture.


\paragraph{Predicting held-out behavior.}
\label{sec:exp:behavioral}s
As a secondary endpoint, we next test whether weight distance predicts behavioral similarity among the models in our controlled experiments. For our translation task, we obtain one deterministic greedy response per prompt and model, embed each output with a small sentence-transformer (\texttt{all-MiniLM-L6-v2}, $384$-d), compute the cosine similarity between models' responses to the same prompt, and average these similarities across prompts. We apply the same procedure to HellaSwag, ARC-C, MMLU, and TruthfulQA, comparing the semantic similarity of each pair of models' responses to the same prompts. We then report the random-effects DerSimonian--Laird \citep{dersimonianlaird1986} pooled correlation $r_{\mathrm{DL}}$ between weight distance and output-embedding similarity, with both quantities $z$-scored within run (placing all models and architectures on one scale). Other generation details are reported in App.~\ref{app:details}.




\paragraph{Comparison to behavior-only estimation.}
\label{sec:exp:blackbox}
We compare our weight-based phylogeny estimation with PhyloLM~\citep{yax2025phylolm}, which estimates model phylogenies from behavior alone and therefore does not require access to model weights. Because our documented HuggingFace family has a known topology, we can score behavior-only estimation against its documented reference lineage. On that family we run the behavior-based PhyloLM method~\citep{yax2025phylolm}, using its released genes and genetic distance on the same checkpoints, and compare it against our weight-based estimation on the identical documented lineage (details in App.~\ref{app:blackbox}). Our controlled comparison runs PhyloLM on every tree for the Llama rank-8 LoRA, rank-64 LoRA, and full-FT conditions (50 reconstructions per condition). We compare on trees with at least one nontrivial unrooted truth split. The encoder--decoder Flan groups are outside PhyloLM's decoder-only scope.


% -----------------------------------------------------------------------
\section{Results}
\label{sec:results}
% -----------------------------------------------------------------------

%\subsection{Weights contain both direct ancestry and phylogenetic signal}

We examine the results of our controlled and real data experiments with our encoder-decoder and decoder models, LoRA trainings, and different task fine-tunings and data sets.


\paragraph{Shared lineage.}
We first determine how well shared lineage estimation (e.g., parent--child or within-branch) performs under different conditions. There are $n_{ord}=26$ eligible trees for this analysis that contain at least one pair of leaves from the same root branch and at least one pair from different root branches. Table~\ref{tab:variants} reports the Mann--Whitney rank-biserial effect and the within-run branch correlation $r_{\mathrm{branch}}$ between cosine distance and a same-branch indicator. Every mean rank-biserial effect is positive and every $r_{\mathrm{branch}}$ is negative. This implies that ordering in a tree has a strong signal in the weights in our experiments.

\begin{table}[h]
\caption{Lineage and recovery by training regime (46 trees). Topology: BHV Fr\'echet mean of
  per-layer NJ trees; pairwise statistics: mean active-layer distance matrix. PAER: all true clades
  recovered ($\mathrm{FN}=0$). Flan is outside PhyloLM's scope. $^\dagger$Cumulative update-space
  all-projections ablation. $^\star$At ceiling: 117/117 truth splits, 46/46 trees.
  $^\ddagger$Distinct cohorts whose statistics coincide at display precision (App.~\ref{app:details}).
  Chance level for a random binary tree: 8.7\% clade recovery, 3.2\% PAER (App.~\ref{app:details}).}
\label{tab:variants}
\begin{center}
\scriptsize
\setlength{\tabcolsep}{1.1pt}
\resizebox{\linewidth}{!}{%
\begin{tabular}{lrrrrrr}
\toprule
Variant & \shortstack{Rank-bis.\\(SE)} & \shortstack{$r_{\mathrm{branch}}$\\(SE)} & \shortstack{Weight clade\\(SE)} & \shortstack{PhyloLM clade\\(SE)} & \shortstack{Weight PAER\\(SE)} & \shortstack{Weight RF (SE);\\FN (SE)} \\
\midrule
\multicolumn{7}{l}{\textbf{Encoder--decoder: Flan-T5-base} ($d=768$)} \\
\multicolumn{7}{l}{\quad\textit{LoRA adapter scope}} \\
\quad LoRA: key only                    & 0.52 (0.09) & $-0.53$ (0.07) & \textbf{92\%} (3\%) & n/a & \textbf{85\%} (5\%) & 2.39 (0.25); 0.17 (0.06) \\
\quad LoRA: q/k/v$^\ddagger$            & 0.56 (0.09) & $-0.62$ (0.07) & \textbf{99\%} (1\%) & n/a & \textbf{98\%} (2\%) & 2.09 (0.21); 0.02 (0.02) \\
\quad LoRA: q/k/v/o (all attention)$^\ddagger$ & 0.56 (0.09) & $-0.62$ (0.07) & \textbf{99\%} (1\%) & n/a & \textbf{98\%} (2\%) & 2.09 (0.21); 0.02 (0.02) \\
\quad LoRA: q/k/v/o + FFN projections$^\dagger$ & 1.00 (0.00) & $-0.91$ (0.02) & \textbf{96\%} (2\%) & n/a & \textbf{89\%} (5\%) & 2.26 (0.21); 0.11 (0.05) \\
\multicolumn{7}{l}{\quad\textit{Full fine-tuning}} \\
\quad Full fine-tuning                  & 0.69 (0.08) & $-0.71$ (0.06) & \textbf{100\%}$^\star$ & n/a & \textbf{100\%}$^\star$ & 2.04 (0.19); 0.00 (0.00) \\
\midrule
\multicolumn{7}{l}{\textbf{Decoder-only: Llama-3.2-1B} ($d=2048$)} \\
\multicolumn{7}{l}{\quad\textit{Controlled exp. trees}} \\
\quad LoRA qkv ($r{=}8$)                 & 0.63 (0.08) & $-0.64$ (0.07) & \textbf{100\%}$^\star$ & 35\% (5\%) & \textbf{100\%}$^\star$ & 2.04 (0.19); 0.00 (0.00) \\
\quad LoRA qkv ($r{=}64$)                & 0.65 (0.08) & $-0.67$ (0.07) & \textbf{100\%}$^\star$ & 24\% (5\%) & \textbf{100\%}$^\star$ & 2.04 (0.19); 0.00 (0.00) \\
\quad Full fine-tuning                   & 0.64 (0.08) & $-0.66$ (0.07) & \textbf{100\%}$^\star$ & 52\% (5\%) & \textbf{100\%}$^\star$ & 1.83 (0.17); 0.00 (0.00) \\
\bottomrule
\end{tabular}
}
\end{center}
\end{table}



\paragraph{Phylogeny estimation.}
\label{sec:results:weighttree}
In addition to ordering, the BHV Fr\'echet means estimate the phylogeny, as shown in Table~\ref{tab:variants}, with respect to clade recovery, polytomy-aware exact recovery (PAER; $\mathrm{FN}=0$), and RF and FN for all eligible phylogenies. In the controlled experiments, clade recovery ranges from 92--100\% and PAER from 85--100\%. The merged LoRA variants recover 92--99\% of clades for Flan and 100\% for Llama; full fine-tuning recovers 100\% and 100\%, respectively. On the same 46 trees with a nontrivial truth split, weights recover 100\% versus PhyloLM's 35\% under rank-8 LoRA, 100\% versus 24\% at rank 64, and 100\% versus 52\% under full fine-tuning.

Treating pairwise layer distance as gene-like estimators reveals that a single attention key matrix can estimate the phylogeny as well as the all-layer aggregate in the full-fine-tuning regime: on the same 46 matched-topology runs, each of the 12 Flan encoder self-attention key matrices on its own (0.24\% of the scalar weights) recovers 100\% of clades with 100\% PAER, as do 10 of the 16 Llama key matrices (0.085\% each), matching the all-weight estimate (App.~\ref{app:layertrace}; Table~\ref{tab:layer_subsets} reports the corresponding layer subsets). The signal also survives low-rank training: q/k/v LoRA adapters that train only 0.54\% of Flan's parameters yield 99\% clade recovery and 98\% PAER, and for Llama at rank 8 the merged q/k/v matrices that carry the update (8.1\% of the weights) yield 100\% and 100\% (Table~\ref{tab:variants}). The individual-layer profile in App.~\ref{app:layertrace} shows that some of the attention layers are strong carriers of lineage signal, although this is architecture dependent.

Finally, we ask whether the diagnostics of \S\ref{sec:framework:theory} predict when estimation succeeds. Figure~\ref{fig:coherence_recovery} plots clade recovery against the fitted Atteson diagnostic $\widetilde m$ and the label-free additivity $A$ for every weight- and behavior-based estimator on the same 46 trees. Both diagnostics predict success. Recovery rises steeply as the distances become more tree-like and then levels off near 100\%. We fit a fractional-logit model to clade estimation $p_c$, the fraction of true clades recovered, with $\operatorname{logit}(p_c)=\beta_0+\beta_1x$. The odds of recovering a clade grow 5.9-fold per $0.10$ increase in $A$ (95\% CI 4.7--8.4, paired bootstrap over trees) and 15.7-fold per tenfold increase in $\widetilde m$ (95\% CI 9.7--58.5). Every weight-based estimator has $A\ge0.82$ and recovers at least 84\% of clades, and the whole-weight runs and layer subsets, all of which have $\widetilde m>1$, recover at least 90\%. In contrast, PhyloLM's behavioral distances have $A=0.58$--$0.63$ and $\widetilde m<0.25$, and they recover 24--52\% of clades. Because $A$ needs no labels, it can indicate how far to trust an estimated phylogeny when the true tree is unknown.

\begin{figure}[t]
  \begin{center}
    \includegraphics[width=\linewidth]{\figpath/fig2_coherence_atteson_v6.pdf}
  \end{center}
  \caption{Mean true-clade recovery against
    \textbf{(A)} the fitted Atteson diagnostic $\widetilde m$ on a logarithmic x-axis and
    \textbf{(B)} label-free quartet
    additivity $A$. 
    Every mark is an average over the same 46 topology-eligible trees used in
    Table~\ref{tab:variants}; vertical intervals on simulation, PhyloLM, and
    subset summaries are $\pm1$ SE. Dashed curves are fractional-logit fits. Odds ratios have
    95\% paired-bootstrap CIs over the 46 trees.}
  \label{fig:coherence_recovery}
\end{figure}


\paragraph{Predicting held-out behavior.}
\label{sec:results:behavioral}

While estimating phylogeny from weights is useful for safety and provenance, it can be of even more use if similarity in weights signals similarity in behavior. Inspired by \citet{cloud2025subliminal}'s subliminal-learning experiment, we test
whether models closer in weight space produce more similar outputs on \textit{held-out} tasks, using translation
(on a new set of trees that never trained on it) and four test tasks that are not in the fine-tuning data (HellaSwag, ARC-Challenge, MMLU, and TruthfulQA).
Because $r_{\mathrm{DL}}$ correlates weight \emph{distance} with output \emph{similarity}, a negative value means that models closer in weight space give more similar outputs, and a positive value means the reverse.
Table~\ref{tab:behavior_holdout} shows a significant negative correlation on every held-out task in Llama. Empty, immediate
end-of-sequence (EOS) responses are observed output and are represented as an explicit empty-output
semantic state in the primary analysis. In the Llama full-FT runs,
38.64\% of HellaSwag leaf outputs were empty, versus 0.02\% for TruthfulQA and 0.00\% for ARC-C
and MMLU. For HellaSwag, including immediate EOS gives $r_{\mathrm{DL}}=-0.17$ [$-0.24,-0.10$], whereas excluding empty completions strengthens the all-node estimate to $-0.27$ [$-0.34,-0.19$] and leaves the other task estimates essentially unchanged (App.~\ref{app:empty-sensitivity}). Thus shared fine-tuning ancestry is associated with more similar behavior on held-out tasks. The effects are modest ($|r_{\mathrm{DL}}|\approx0.2$) but significant.
Flan has inconsistent results. Its HellaSwag correlation is significantly positive, and its held-out translation correlation is near zero. Part of this reflects degenerate outputs after repeated Flan fine-tuning, where 15--30\% of Flan outputs are only a label from a fine-tuning task, such as \texttt{entailment}. The Flan row therefore measures similarity among outputs rather than task performance (App.~\ref{app:regression}).



\begin{table}[t]
\caption{Weight$\to$behavior bridge under full fine-tuning: random-effects pooled
  $r_{\mathrm{DL}}$ with 95\% CI. Immediate EOS is represented explicitly. Flan output-form
  contamination and heterogeneity are detailed in App.~\ref{app:regression}.}
\label{tab:behavior_holdout}
\centering
\small
\setlength{\tabcolsep}{5pt}
\begin{tabular}{lcccccc}
\toprule
$r_{\mathrm{DL}}$ [95\% CI] & \multicolumn{2}{c}{Translation} & HellaSwag & ARC-C & MMLU & TruthfulQA \\
\cmidrule(lr){2-3}
Model & trained & held-out & held-out & held-out & held-out & held-out \\
\midrule
Llama-3.2-1B & \shortstack{$-0.20$\\{\scriptsize$[-.26,-.13]$}} & \shortstack{$-0.23$\\{\scriptsize$[-.29,-.17]$}} & \shortstack{$-0.17$\\{\scriptsize$[-.24,-.10]$}} & \shortstack{$-0.20$\\{\scriptsize$[-.26,-.13]$}} & \shortstack{$-0.19$\\{\scriptsize$[-.26,-.12]$}} & \shortstack{$-0.22$\\{\scriptsize$[-.26,-.16]$}} \\
\addlinespace
Flan-T5-base & \shortstack{$-0.10$\\{\scriptsize$[-.16,-.04]$}} & \shortstack{$+0.04$\\{\scriptsize$[-.01,+.09]$}} & \shortstack{$+0.09$\\{\scriptsize$[+.05,+.14]$}} & \shortstack{$-0.14$\\{\scriptsize$[-.19,-.09]$}} & \shortstack{$-0.16$\\{\scriptsize$[-.20,-.11]$}} & \shortstack{$-0.17$\\{\scriptsize$[-.22,-.12]$}} \\
\bottomrule
\end{tabular}
\end{table}


\paragraph{Weight- versus behavior-based estimation.}
\label{sec:results:blackbox}

When weights are unavailable or architectures differ, lineage must be inferred from behavior alone. Table~\ref{tab:variants} already compares our estimates with the behavior-only PhyloLM~\citep{yax2025phylolm} on the controlled trees. We also score both approaches against documented HuggingFace families. On three such families (Mistral-7B, Llama-2-7B, and Flan-T5-base), our weight-based estimates recover every documented split (App.~\ref{app:organic-additional}). On the ten-model Mistral family, PhyloLM also recovers all three forks, although the weight-based distances are more additive (informative-quartet $A=0.98$ vs.\ $0.83$, Fig.~\ref{fig:real-trees}).

The Mistral forks are large RLHF or direct-preference-optimization (DPO) updates that shift behavior substantially. Our controlled fine-tunings are smaller, and there PhyloLM recovers only 24--52\% of clades versus 100\% for weights (Table~\ref{tab:variants}). Behavior thus carries some, but not all, of the lineage signal in the weights, and weight-based estimation does not depend on a choice of prompts. Overall, weight-based estimation matches behavior-based estimation on this real family and substantially outperforms it in the controlled experiments.

% -----------------------------------------------------------------------
\section{Discussion}
\label{sec:discussion}
% -----------------------------------------------------------------------

We show that model family phylogeny can be estimated from weights alone, and that using a single attention key matrix---0.1--0.2\% of the scalar weights---estimates training trees as well as using all analyzed layers. We also find that clade estimation works across encoder--decoder and decoder-only models and survives LoRA training at low rank ($r=8$). Across our controlled conditions, PAER ($\mathrm{FN}=0$) ranges from 85--100\%. The label-free additivity proxy $A$ is associated with increased clade recovery in our experiments, although the relationship depends on architecture and layer subset. In addition to estimating phylogeny and shared lineage from weights, we show that models closer in weight space tend to have more similar outputs on held-out tasks. The effect depends on the task and ranges from $r_{\mathrm{DL}}=-0.23$ to $-0.17$ across the held-out Llama tasks. Finally, weight- and behavior-based estimation recover the same documented Mistral family, while weight-based estimation is substantially stronger on several controlled trees and requires no prompts.


\paragraph{Implications.} Our results suggest three uses beyond reconstructing controlled trees. First, weight-based phylogenies can fingerprint models and help chart the largely undocumented atlas of fine-tuned models on public hubs \citep{horwitz2025chartatlasworldsmodels}, at least within a shared architecture. Second, four-point additivity offers a label-free confidence check. Because clade recovery rises sharply with $A$ (odds ratio 5.9 per $0.10$), a practitioner without a ground-truth tree can use $A$ to judge how far to trust an estimated lineage, which brings a classical idea from phylogenetics to model auditing. Third, a tiny, architecture-dependent fraction of the weights suffices. A single attention key matrix, 0.1--0.2\% of the weights, matches the all-weight estimate, which makes lineage audits cheap even for large models.

\paragraph{Limitations.} One limitation is that our methods for phylogeny estimation using weights currently only work when the underlying architecture is identical, i.e., the weight space is the same.  This means that Llama cannot directly be compared to Mistral, for example. Moreover, our method assumes that all models are leaves unlike \cite{horwitz2025mother}. Another issue is that our controlled training tree generation results from a fairly simple process that may not result in the most realistic training trees.  We attempt to mitigate that with three documented real-world families. Future directions include better relating models across different architectures such as with CKA and further exploring the mechanism by which lineage is inherited. In summary, we hope that our work continues to illuminate the need to consider models holistically rather than only as a collection of benchmark scores, and shows how analyzing weight space can shed light on model behavior, especially as traits passed through fine-tuning.


% -----------------------------------------------------------------------
% Author acknowledgments omitted for double-blind review; add at camera-ready.

\subsection*{Reproducibility statement}
All controlled experiments follow seeded, fully specified procedures. Appendix~\ref{app:details}
describes the tree-generation process (Poisson branching, node budget, depth limit, and acceptance
rule), the 36 task/dataset pairs and their assignment, the training configuration of both base models
(pinned HuggingFace revisions, learning rates, batching, schedules, precision, sequence limits, and
LoRA hyperparameters; Table~\ref{tab:training_config}), and the behavioral test-task protocol (prompt
counts, decoding settings, generation seeds, and the pinned sentence-embedding model).
Appendix~\ref{app:layers} defines every layer subset and LoRA variant (Tables~\ref{tab:layer_subsets}
and~\ref{tab:lora_variants}); Appendix~\ref{app:blackbox} gives the PhyloLM configuration and pinned
commit used for the behavior-only comparison; and Appendix~\ref{app:organic-additional} names the
checkpoints in the three documented HuggingFace families. The scoring and statistical procedures
(polytomy-aware RF/FN/PAER scoring, the BHV Fr\'echet mean, the four-point and fitted Atteson
diagnostics, and jackknife variances pooled by DerSimonian--Laird) are specified in
Sections~\ref{sec:framework}--\ref{sec:experiments} and Appendix~\ref{app:regression}. Code for tree
generation, training, distance computation, reconstruction, scoring, and analysis, together with the
tree manifests, ground-truth trees, and per-tree distance matrices needed to regenerate every table
and figure, will be released publicly upon publication.
% Submission TODO: add the anonymous repository link to the sentence above once the blind snapshot exists.

\subsection*{AI use statement}
Generative AI assistants were used to write and debug software across the experimental
pipeline (data processing, weight-distance computation, tree reconstruction and scoring, statistical
analysis, figures, and tables) and to review and edit the manuscript, including checking consistency between numerical results and prose,
improving readability, searching and summarizing relevant literature, and formatting citations.  The research questions, experimental design,
hypotheses, conceptual framework, and scientific interpretations are the authors' own. The authors
reviewed and tested AI-assisted code and independently verified all results, figures, numerical claims,
citations, and manuscript text. The authors take responsibility for the final content of this work,
including all text, claims, and artifacts produced with the aid of generative AI.

% -----------------------------------------------------------------------
\bibliography{references}
\bibliographystyle{iclr2027_conference}


% -----------------------------------------------------------------------
\appendix
% -----------------------------------------------------------------------

\section{Experimental Details}
\label{app:details}

\subsection*{Tree Generation}
\label{app:tree-gen}

Training trees are generated via a Poisson branching process as follows:

\begin{enumerate}
    \item The virtual root draws $k \sim \mathrm{Poisson}(\lambda = 1.5)$ children, with a minimum of 1. We instantiate the same accepted topologies and task/data assignments with either Flan-T5-base or Llama-3.2-1B.
    \item Each queued node below the depth limit draws $k \sim \mathrm{Poisson}(\lambda = 1.5)$ children while node budget remains. Nodes that draw $k = 0$ naturally become leaves; no separate pruning step is applied.
    \item Expansion is breadth-first, with a budget of at most $n_\mathrm{nodes} = 14$ training nodes (not counting the root) and maximum depth $n_\mathrm{layers} = 4$. Reaching depth 4 on one branch prevents only those depth-4 nodes from expanding; shallower nodes already in the queue can still produce children while budget remains. Growth ends when the queue is empty or the node budget is exhausted. Candidate trees with fewer than four leaves are rejected, and seeded draws continue until 50 trees are accepted.
    \item Each non-root node is assigned one of 36 task/dataset pairs (9 summarization, 10 classification, 8 QA, and 9 translation). The pairs are sampled without replacement within a tree, so datasets do not repeat but task families can. Topology generation and task/dataset assignment are seeded separately.
    \item Topology recovery uses only accepted trees with at least one nontrivial unrooted truth split---an internal edge that partitions the leaves into two groups of at least two. Four accepted trees each have four leaves but reduce to unresolved stars, with only trivial one-versus-three splits; this leaves 46 trees for recovery. Four leaves alone is not exclusionary.
\end{enumerate}

This procedure produces trees with 4--10 observed leaves and variable depth and branching
structure. In our recovery experiments, we use the 46 trees that have at least one nontrivial
unrooted truth split. Across these 46 trees, there are 609 trained nodes (349 leaves and 260 internal
nodes), in addition to 46 untrained roots. The ground-truth topology is fully known by construction,
enabling exact RF evaluation.

\paragraph{Chance baseline.} The chance level quoted in the caption of Table~\ref{tab:variants} is an
exact expectation rather than a simulation. There are $(2n-5)!!$ unrooted binary topologies on $n$
labelled leaves \citep{sempleSteel2003phylogenetics}. A fixed split whose sides contain $a$ and
$b=n-a$ leaves is present in $(2a-3)!!\,(2b-3)!!$ of them, because contracting the split edge
leaves a rooted binary tree on each side and there are $(2k-3)!!$ rooted binary trees on $k$ leaves.
A uniformly random binary tree therefore contains the split with probability
$(2a-3)!!\,(2b-3)!!/(2n-5)!!$: one third for a four-leaf quartet, and $0.5\%$ for a balanced
$5\,|\,5$ split of ten leaves. Averaging this probability over the truth splits of each of the 46
topology-eligible trees, and then over trees exactly as clade recovery is averaged in
Table~\ref{tab:variants}, gives an expected chance clade recovery of $8.7\%$ (SE $1.1$ across
trees). A random tree recovers every split of a truth tree $T$ (PAER) only if it refines $T$; a
multifurcating $T$ has $\prod_v (2d_v-5)!!$ binary refinements, the product running over its
internal vertices $v$ of degree $d_v$, so the chance PAER is the mean over trees of
$\prod_v (2d_v-5)!!/(2n-5)!!$, which is $3.2\%$.

\paragraph{Experiment structure.} All topology-recovery comparisons use the same 46 retained
unlabeled tree topologies, with node identities aligned across conditions. In the ordinary Flan
and Llama cohorts, every aligned node also retains the same task/dataset label across root models
and fine-tuning interventions. The matched no-translation cohorts preserve those topologies
and node identities, but replace translation assignments with unused classification, QA, or
summarization datasets. Thus, comparisons are paired by tree and,
except for the intended translation-removal intervention, by node label.

Different analyses use fixed eligible subsets of these shared trees. Topology recovery uses the
common 46 trees with at least one nontrivial unrooted truth split (excluding topology IDs 015, 020,
037, and 047), while direct ordering uses the common 26 trees with a defined cross-root-branch
comparison. PhyloLM produces a behavioral reconstruction for every tree in each matched Llama
condition, and its topology summaries and the matched weight values in
Table~\ref{tab:variants} use the same 46-tree recovery subset. The Flan q/k/v and q/k/v/o LoRA cohorts are
trained and analyzed separately (separate adapters, checkpoints, and distance matrices); their
estimated trees differ only by a polytomy-resolving split on five of the 50 trees, so their
Table~\ref{tab:variants} rows coincide at display precision. The Flan and Llama per-layer profiles
in Fig.~\ref{fig:layertrace_grid} likewise use their matched 46-tree full-fine-tuning cohorts.

\subsection*{Training Configuration}
Table~\ref{tab:training_config} summarizes the audited configurations for
\texttt{google/flan-t5-base}\footnote{Pinned revision:
\nolinkurl{7bcac572ce56db69c1ea7c8af255c5d7c9672fc2}.} and the base, non-Instruct
\texttt{meta-llama/Llama-3.2-1B}\footnote{Pinned revision:
\nolinkurl{4e20de362430cd3b72f300e6b0f18e50e7166e08}.} Each child is initialized
from its parent and trained under the configuration for its cohort.

\begin{center}
\begin{minipage}{\linewidth}
\centering
\refstepcounter{table}
\label{tab:training_config}
\textbf{Table \thetable:} Training configuration for the controlled Flan and Llama cohorts.
\vspace{4pt}

\small
\setlength{\tabcolsep}{4pt}
\begin{tabular}{@{}p{0.19\linewidth}p{0.38\linewidth}p{0.38\linewidth}@{}}
\toprule
Setting & Flan-T5-base & Llama-3.2-1B \\
\midrule
Objective & Sequence-to-sequence & Completion-only causal language modeling \\
Training length & 2,000 steps; no early stopping & 2,000 steps; no early stopping \\
Learning rate & $3\times10^{-4}$ (full FT and LoRA) & $2\times10^{-5}$ (full FT); $3\times10^{-4}$ (LoRA) \\
Batching & Batch size 8; gradient accumulation 4; effective batch 32 & Batch size 8; gradient accumulation 4; effective batch 32 \\
Schedule & 200 warmup steps; weight decay 0.01 & 200 warmup steps; weight decay 0.01 \\
Precision & bfloat16 & bfloat16 \\
Sequence limits & 512 input tokens; 128 target tokens & 1,024 tokens total for prompt plus completion \\
Padding & Right & Right \\
LoRA & $r=8$, $\alpha=16$, dropout 0.05; module scopes in Table~\ref{tab:lora_variants} & QKV; $(r,\alpha)=(8,16)$ or $(64,128)$; dropout 0.05 \\
Data sample & At most 10,000 train and 1,000 evaluation examples per dataset & At most 10,000 train and 1,000 evaluation examples per dataset \\
\bottomrule
\end{tabular}
\end{minipage}
\end{center}

Training and dataset subsampling are seeded and held fixed across matched conditions and across
nodes assigned the same dataset. Topology generation and task/dataset assignment use separate seed
streams, as described above.

\subsection*{Behavioral Test Tasks and Analysis Population}
Within each behavioral cohort, the primary analysis uses all 641 trained, non-root checkpoints
across 50 trees: 276 trained internal nodes and 365 terminal leaves. Each internal node is an actual
checkpoint saved during training, not a latent ancestor inferred from its descendants; we evaluate
its weights and outputs directly, just as we do for a leaf. For each within-tree checkpoint pair, weight distance is the mean
active-layer cosine distance described above. Translation uses 200 frozen prompts with one
deterministic greedy completion per checkpoint (up to 128 new tokens for Flan and 64 for Llama).
HellaSwag, ARC-Challenge, MMLU, and TruthfulQA
each use 100 frozen prompts and three sampled completions per prompt, with a distinct generation seed
for each draw (temperature 1, $p=1$; up to 64 new tokens). Dolly, used only in the
end-of-sequence (EOS) sensitivity, uses 30 prompts, three draws, and up to 96 new tokens.

Outputs are embedded with \texttt{all-MiniLM-L6-v2}\footnote{Pinned revision:
\nolinkurl{1110a243fdf4706b3f48f1d95db1a4f5529b4d41}.}. We average prompt-paired cosine similarity (and
prompt-by-draw similarity for the three-draw tasks) before correlating it with weight distance within
each tree. The inference unit is the tree. Within-tree Fisher-$z$ variances are estimated by deleting one
trained node at a time, after which the tree correlations are pooled by
DerSimonian--Laird~\citep{dersimonianlaird1986}. Pair counts are descriptive and are not treated as
independent observations.

\section{Layer Variant Definitions}
\label{app:layers}
Two families of variants appear in the paper. \emph{layer subsets} (Table~\ref{tab:layer_subsets})
restrict the full-weight distance to a chosen set of layers, testing which layers carry the lineage
signal. For each subset, we average the selected per-layer cosine-distance matrices before neighbor
joining. ``SA.k'' denotes the self-attention key projection, and ``pre-attention LayerNorm'' denotes
\texttt{layer.0.layer\_norm} in Flan-T5-base and \texttt{input\_layernorm} in Llama-3.2-1B.
Counts and percentages are relative to the 282 Flan layers or 146 Llama layers in the analysis.
The \texttt{low\_signal} control contains post-attention, cross-attention, and post-FFN LayerNorms
for Flan and post-attention LayerNorms for Llama.

\emph{LoRA variants} (Table~\ref{tab:lora_variants}) instead change which modules are fine-tuned.
All use dropout $0.05$; the Flan variants and Llama rank-8 variant use $\alpha=16$, while the Llama
rank-64 variant uses $\alpha=128$. Their controlled-tree analysis merges the full
root-to-node adapter chain so distances are computed between standalone full-weight leaves. The
all-projections row is retained as an explicitly marked cumulative-update ablation: it compares the
cumulative root-to-node adapter displacement with the shared base omitted, whereas the other LoRA rows
use standalone merged full weights. It is therefore excluded from the merged-weight diagnostic fit.
In Table~\ref{tab:lora_variants}, $q,k,v,o$ denote the attention projections; $\mathrm{wi}_{0,1}$
and $\mathrm{wo}$ denote the two gated FFN inputs and the FFN output.

\begin{table}[h]
\caption{Full-FT layer-subset recovery on the matched 46-tree cohorts (cosine). Values are
percentages reported as mean (SE); PAER is polytomy-aware exact recovery.}
\label{tab:layer_subsets}
\centering
\scriptsize
\setlength{\tabcolsep}{2pt}
\begin{tabular}{lllrrrr}
\toprule
Model & Variant & Layers included & Count & \% layers & Clade (SE) & PAER (SE) \\
\midrule
Flan & \texttt{full}         & all layers                                  & 282 & 100.0 & 100.0 (0.0) & 100.0 (0.0) \\
     & \texttt{high\_signal} & SA.k $+$ pre-attn LN                        & 48  & 17.0  & 100.0 (0.0) & 100.0 (0.0) \\
     & \texttt{sak}          & SA.k only                                    & 24  & 8.5   & 100.0 (0.0) & 100.0 (0.0) \\
     & \texttt{encoder\_h}   & encoder SA.k $+$ pre-attn LN                 & 24  & 8.5   & 100.0 (0.0) & 100.0 (0.0) \\
     & \texttt{enc\_sak}     & encoder SA.k only                            & 12  & 4.3   & 100.0 (0.0) & 100.0 (0.0) \\
     & \texttt{low\_signal}  & LayerNorm control                            & 36  & 12.8  & 100.0 (0.0) & 100.0 (0.0) \\
\midrule
Llama & \texttt{full}         & all layers                                  & 146 & 100.0 & 100.0 (0.0) & 100.0 (0.0) \\
      & \texttt{high\_signal} & SA.k $+$ pre-attn LN                        & 32  & 21.9  & 100.0 (0.0) & 100.0 (0.0) \\
      & \texttt{sak}          & SA.k only                                    & 16  & 11.0  & 100.0 (0.0) & 100.0 (0.0) \\
      & \texttt{low\_signal}  & LayerNorm control                            & 16  & 11.0  & 53.3 (5.0)  & 21.7 (6.1) \\
\bottomrule
\end{tabular}
\end{table}

The zero SEs are exact rather than rounded: every Flan subset and each of the three signal-bearing
Llama subsets recovers all 46 eligible topologies. The Llama LayerNorm control does not, showing that
the ceiling result is specific to the informative subsets rather than an artifact of the evaluation.

\begin{table}[h]
\caption{LoRA configurations used in the fixed-2,000-step cohorts. Percentages compare trainable
LoRA parameters with the analyzed base-model weights.}
\label{tab:lora_variants}
\centering
\scriptsize
\begin{tabular}{lllrrrr}
\toprule
Model & Variant & Target modules & $r$ & $\alpha$ & Trainable params & \% weights \\
\midrule
Flan & \texttt{lora\_k}         & $k$                                      & 8  & 16  & 442,368   & 0.18 \\
     & \texttt{lora\_qkv}       & $q,k,v$                                  & 8  & 16  & 1,327,104 & 0.54 \\
     & \texttt{lora\_qkvo}      & $q,k,v,o$                                & 8  & 16  & 1,769,472 & 0.71 \\
     & \texttt{lora\_all\_proj} & $q,k,v,o,\mathrm{wi}_{0,1},\mathrm{wo}$ & 8  & 16  & 3,391,488 & 1.37 \\
\midrule
Llama & \texttt{lora\_qkv}      & $q,k,v$                                  & 8  & 16  & 1,179,648 & 0.10 \\
      & \texttt{lora\_qkv}      & $q,k,v$                                  & 64 & 128 & 9,437,184 & 0.76 \\
\bottomrule
\end{tabular}
\end{table}

\subsection*{Per-layer Recovery}
\label{app:layertrace}
Treating each layer as an independent estimator, Fig.~\ref{fig:layertrace_grid} reports the per-layer
FN-rate across architectures. On the Flan trees the signal is diffuse across individual layers:
every encoder attention and feed-forward matrix, including each self-attention key matrix, recovers all true clades in
all 46 runs, while decoder layers miss between 0\% and 12\% on average. In Llama, 10 of the 16
self-attention key matrices likewise recover all clades in all runs, and the remaining trained
projections miss about 1\%. This is consistent with Table~\ref{tab:layer_subsets}, where averaging
only the 12 encoder \textsc{SelfAttention.k} matrices matches the all-weight reconstruction. \textsc{LayerNorm} produces
prominent local maxima in both architectures.

\begin{figure}[h]
  \begin{center}
    \includegraphics[width=\linewidth]{\figpath/fig3_layertrace_grid.pdf}
  \end{center}
  \caption{Block-aligned per-layer FN-rate (\% of true clades missed; mean $\pm$ 2\,SE,
    cosine, full fine-tuning) for the encoder and decoder of Flan-T5-base ($d{=}768$, $n=46$)
    and the decoder of the matched Llama-3.2-1B cohort ($d{=}2048$, $n=46$). Each transformer block
    occupies one x-position; gray points show the remaining layers without connecting unrelated
    layer families. The two Flan block-0
    relative-attention-bias layers (32.8\% encoder; 55.4\% decoder) are omitted from the plot.}
  \label{fig:layertrace_grid}
\end{figure}

\section{Blackbox Experiment Details}
\label{app:blackbox}

\paragraph{PhyloLM reproduction.} For the behavior-only comparison on the documented Mistral family
(Fig.~\ref{fig:real-trees}), we reproduce PhyloLM \citep{yax2025phylolm} faithfully using its released
\texttt{math} gene bank\footnote{Pinned PhyloLM revision:
\nolinkurl{8c70edf062a0adce2a3e6c8c79cd23a645fd0905}.} and its published Nei-similarity genetic
distance. We sample $G=128$ genes and, for each, draw $N=32$ raw completions of $4$ tokens per model
(temperature $1$, no chat template, matching PhyloLM's design), take the first $4$ characters as the
allele, form the per-gene allele populations, compute the Nei similarity $S$ between every pair of the
ten checkpoints, convert it to $-\log\max(S,10^{-3})$, and build the tree by neighbor-joining.
The reconstruction is scored against the same documented truth tree and with the same polytomy-aware
scorer used for the whitebox families.

The controlled comparison uses the same locked PhyloLM contract ($G=128$, $N=32$, four-token
continuations, temperature $1$) with seeded sampling on the 50 matched-topology Llama
rank-8 LoRA, rank-64 LoRA, and full-FT conditions. For LoRA leaves, we reconstruct and query the
ordered cumulative root-to-leaf adapter chain; full-FT leaves are loaded directly. Each condition
contains generated populations for all 365 leaves across the 50 trees, and all generated alleles are
nonempty. The four trees whose truths contain no nontrivial unrooted split are valid behavioral
reconstructions but cannot define clade recovery, PAER, RF, or FN, leaving the same $n=46$ topology
sample in every row of Table~\ref{tab:variants}.

\section{Additional Organic Lineage Checks}
\label{app:organic-additional}

The paper evaluates three documented organic HuggingFace lineages in total. The ten-checkpoint
Mistral-7B family contains three documented nontrivial splits and supports the matched PhyloLM
comparison (Fig.~\ref{fig:real-trees}). Weights and PhyloLM both recover all three splits. Two smaller
checks establish that the weight-based result is not unique to Mistral. In a seven-checkpoint Flan-T5-base lineage, averaged-matrix NJ and the BHV Fr\'echet
mean recover the sole documented Flan-Alpaca descendant clade (FN $=0/1$); its complementary split has
98.2\% Fr\'echet support under cosine distance. In a ten-checkpoint Llama-2-7B lineage, the cosine
Fr\'echet estimate likewise recovers the sole documented seven-model chat-descendant clade (FN
$=0/1$), with 95.2\% support for the complementary base/Nous-Hermes/Vicuna split. These two references
are bushy and specify no ordering within the recovered clade, so we claim recovery of one documented
split in each, not complete resolution of their descendant histories. Figures~\ref{fig:organic-flan}
and~\ref{fig:organic-llama} show the documented families alongside their reconstructions.

\begin{figure}[h]
  \begin{center}
    \includegraphics[width=\linewidth]{\figpath/hf_mistral_hero_3panel.pdf}
  \end{center}
  \caption{\textbf{Documented Mistral lineage:} truth (left), our weight-based estimate (middle),
    and PhyloLM's output-based estimate (right)~\citep{yax2025phylolm}. Both recover all three splits.}
  \label{fig:real-trees}
\end{figure}

\begin{figure}[h]
  \begin{center}
    \includegraphics[width=\linewidth]{\figpath/hf_flan_trees_vs_truth.png}
  \end{center}
  \caption{Flan-T5-base organic family. The cosine Fr\'echet reconstruction recovers the documented
  Flan-Alpaca split. Relationships not specified by the reference polytomy are not treated as ground
  truth.}
  \label{fig:organic-flan}
\end{figure}

\begin{figure}[h]
  \begin{center}
    \includegraphics[width=\linewidth]{\figpath/hf_llama2_trees_vs_truth.png}
  \end{center}
  \caption{Llama-2-7B organic family. The cosine Fr\'echet reconstruction recovers the documented
  chat-descendant split; additional resolution within that clade is not claimed.}
  \label{fig:organic-llama}
\end{figure}

\section{Behavioral Meta-Analysis Diagnostics}
\label{app:regression}

\paragraph{Pooling estimator.} Because each tree has its own weight-space scale, a raw pooled slope is
not meaningful across trees. We standardize within tree and estimate the variance of each Fisher-$z$
correlation with a delete-one-trained-node jackknife, then pool tree-level effects using a
DerSimonian--Laird random-effects meta-analysis. Pair counts are descriptive and are never treated as
independent observations.

\paragraph{Per-cell heterogeneity.} Table~\ref{tab:behavior_het} gives, for every cell of the main-text
bridge (Table~\ref{tab:behavior_holdout}, paired semantic DV), $r_{\mathrm{DL}}$ with its $95\%$ CI,
Higgins' $I^2$~\citep{higgins2003measuring}, and the number of trees $k$. The statistic estimates the
proportion of observed effect variation attributable to between-tree heterogeneity rather than
sampling error; here it ranges from none to moderate ($I^2$ $0$--$36\%$). All cells except Flan
translation-withheld exclude zero. The Llama cells are negative, Flan HellaSwag is significantly
positive, the other significant Flan cells are negative, and Flan translation-withheld is near zero.

\paragraph{Flan output-form audit.} We re-audited the finalized full-fine-tuning response sets: 641
trained nodes in the ordinary Flan cohort and a separate 641 in the matched translation-withheld
cohort. ARC-C and MMLU prompts display answer choices, so valid option-only answers are not counted as
contamination; HellaSwag is rendered as a free completion and TruthfulQA as a concise free response.
Even after separating valid option labels, outputs that exactly match unrelated fine-tuning-task
labels (e.g., \texttt{entailment}, \texttt{positive}, or \texttt{sci/tech}) comprise 20.4\% of ARC-C
outputs and 22.9\% of MMLU outputs. The same contamination appears throughout the Flan panel: 14.7\%
for trained translation, 22.6\% for held-out translation, 21.3\% for HellaSwag, and 30.3\% for
TruthfulQA. Selective task-level ``degenerate'' flags would therefore be misleading. We report every
cell but interpret the Flan row as associations among generated outputs, not as evidence of task
performance.

\begin{table}[h]
\caption{Per-cell heterogeneity for the weight$\to$behavior bridge (Table~\ref{tab:behavior_holdout};
  paired semantic DV with immediate-EOS completions represented explicitly, full fine-tuning).
  DerSimonian--Laird random-effects $r_{\mathrm{DL}}$ with $95\%$ CI, Higgins $I^2$,
  and number of trees $k$. Bold rows have 95\% CIs excluding zero.}
\label{tab:behavior_het}
\begin{center}
\small
\begin{tabular}{llrrr}
\toprule
Model & Task & $r_{\mathrm{DL}}$ [95\% CI] & $I^2$ & $k$ \\
\midrule
\sigrow{Llama-3.2-1B}{Translation (trained)}{-0.20\ [-0.26,-0.13]}{13\%}{50}
\sigrow{Llama-3.2-1B}{Translation (held-out)}{-0.23\ [-0.29,-0.17]}{0\%}{50}
\sigrow{Llama-3.2-1B}{HellaSwag (held-out)}{-0.17\ [-0.24,-0.10]}{26\%}{50}
\sigrow{Llama-3.2-1B}{ARC-C (held-out)}{-0.20\ [-0.26,-0.13]}{33\%}{50}
\sigrow{Llama-3.2-1B}{MMLU (held-out)}{-0.19\ [-0.26,-0.12]}{36\%}{50}
\sigrow{Llama-3.2-1B}{TruthfulQA (held-out)}{-0.22\ [-0.27,-0.16]}{19\%}{50}
\sigrow{Flan-T5-base}{Translation (trained)}{-0.10\ [-0.16,-0.04]}{0\%}{50}
Flan-T5-base & Translation (withheld)             & $+0.04$ [$-0.01,+0.09$] & 0\% & 50 \\
\sigrow{Flan-T5-base}{HellaSwag (held-out)}{+0.09\ [+0.05,+0.14]}{0\%}{50}
\sigrow{Flan-T5-base}{ARC-C (held-out)}{-0.14\ [-0.19,-0.09]}{0\%}{50}
\sigrow{Flan-T5-base}{MMLU (held-out)}{-0.16\ [-0.20,-0.11]}{0\%}{50}
\sigrow{Flan-T5-base}{TruthfulQA (held-out)}{-0.17\ [-0.22,-0.12]}{3\%}{50}
\bottomrule
\end{tabular}
\end{center}
\end{table}


\section{Immediate End-of-Sequence (EOS) Sensitivity and Depth}
\label{app:empty-sensitivity}

An immediate EOS is an observed model response rather than a missing observation. The primary semantic
analysis therefore represents an empty completion as an explicit semantic state. Table~\ref{tab:eos-sensitivity}
reports the complementary sensitivity analysis that removes pairs involving an empty completion. The choice
matters materially only for HellaSwag, where immediate EOS is common.

\begin{table}[h]
\caption{Llama-3.2-1B full-FT all-trained-node semantic bridge with immediate-EOS outputs included as an explicit
  semantic state (primary) or excluded (sensitivity). Values are $r_{\mathrm{DL}}$ [95\% CI]. Empty
  rates are pooled over all leaf outputs and all three draws.}
\label{tab:eos-sensitivity}
\centering
\small
\begin{tabular}{lrrr}
\toprule
Task & EOS included & EOS excluded & Empty outputs \\
\midrule
ARC-Challenge & $-0.195$ [$-0.257,-0.131$] & $-0.195$ [$-0.257,-0.131$] & $0.00\%$ \\
Dolly open-ended & $-0.081$ [$-0.134,-0.028$] & $-0.083$ [$-0.136,-0.030$] & $0.16\%$ \\
HellaSwag & $-0.175$ [$-0.245,-0.103$] & $-0.265$ [$-0.338,-0.189$] & $38.64\%$ \\
MMLU & $-0.189$ [$-0.256,-0.121$] & $-0.189$ [$-0.256,-0.121$] & $0.00\%$ \\
TruthfulQA & $-0.218$ [$-0.275,-0.161$] & $-0.219$ [$-0.275,-0.161$] & $0.02\%$ \\
\bottomrule
\end{tabular}
\end{table}

Table~\ref{tab:eos-depth} reports the same empty-output rate by leaf depth. HellaSwag empties increase
from 16.72\% at depth 1 to 41.13\% at depth 3, then remain essentially level at 40.25\% at depth 4. Immediate termination is therefore itself
associated with accumulated fine-tuning depth; excluding these outputs nevertheless strengthens the
negative all-node semantic association, as shown in Table~\ref{tab:eos-sensitivity}.

\begin{table}[h]
\caption{Immediate-EOS output rate by leaf depth in the Llama-3.2-1B full-FT evaluation trees. Rates
  pool outputs across all leaves observed at each depth and all three draws.}
\label{tab:eos-depth}
\centering
\small
\begin{tabular}{lrrrr}
\toprule
Task & Depth 1 & Depth 2 & Depth 3 & Depth 4 \\
\midrule
ARC-Challenge & 0.00\% & 0.00\% & 0.00\% & 0.00\% \\
Dolly open-ended & 1.45\% & 0.34\% & 0.10\% & 0.06\% \\
HellaSwag & 16.72\% & 30.86\% & 41.13\% & 40.25\% \\
MMLU & 0.00\% & 0.00\% & 0.00\% & 0.00\% \\
TruthfulQA & 0.10\% & 0.04\% & 0.01\% & $<0.01\%$ \\
\bottomrule
\end{tabular}
\end{table}


\end{document}
